\documentclass[aspectratio=169]{beamer}
\usetheme{Madrid}
\usecolortheme{default}
\usepackage[utf8]{inputenc}
\usepackage[T5]{fontenc}
\usepackage{graphicx}
\usepackage{hyperref}
\usepackage{booktabs}
\usepackage{amsmath}

\setbeamertemplate{caption}[numbered]
\renewcommand{\figurename}{Hình}
\renewcommand{\thefigure}{4.\arabic{figure}}

\title[Chương 4: Phân loại \& Hồi quy]{Học Sâu (Deep Learning) \\ \vspace{0.5cm} \Large Chương 4: Bắt đầu với Mạng Nơ-ron (Phân loại và Hồi quy)}
\author{Đại học Đông Á}
\date{\today}

\begin{document}

\begin{frame}
    \titlepage
\end{frame}

\begin{frame}{Nội dung chương}
    \tableofcontents
\end{frame}

\section{1. Các thuật ngữ cốt lõi (Glossary)}

\begin{frame}{Khái niệm Cơ bản}
    Trong phân loại và hồi quy, các kỹ sư Học máy sử dụng thuật ngữ rất nghiêm ngặt:
    \begin{itemize}
        \item \textbf{Sample (Mẫu):} Một điểm dữ liệu đầu vào. Ví dụ: Một bức ảnh, một đoạn văn bản.
        \item \textbf{Target (Mục tiêu):} Giá trị thực tế (Ground-truth). Ví dụ: Nhãn "Chó" cho ảnh chó.
        \item \textbf{Prediction (Dự đoán):} Giá trị đầu ra từ mô hình.
        \item \textbf{Loss value (Giá trị mất mát):} Khoảng cách đo lường sự chênh lệch giữa Prediction và Target.
    \end{itemize}
\end{frame}

\begin{frame}{Thuật ngữ về Nhãn (Labels \& Classes)}
    \begin{itemize}
        \item \textbf{Classes (Lớp):} Tập hợp các danh mục khả dĩ. (Ví dụ: Chó, Mèo, Chim).
        \item \textbf{Label (Nhãn):} Một phiên bản cụ thể của Class được gán cho một Sample.
        \item \textbf{Ground-truth (Sự thật gốc):} Toàn bộ các targets của tập dữ liệu, thường được con người gán nhãn thủ công.
        \item \textbf{Batch (Mini-batch - Lô):} Một tập hợp nhỏ các samples (thường 8-128) được mô hình xử lý cùng lúc để tính gradient.
    \end{itemize}
\end{frame}

\begin{frame}{Phân loại các bài toán Học máy}
    Có 3 bài toán phổ biến nhất khi mới bắt đầu:
    \begin{enumerate}
        \item \textbf{Phân loại Nhị phân (Binary Classification):} Phân loại đầu vào thành 2 nhóm độc quyền (Ví dụ: Chó hoặc Mèo, Spam hoặc Non-Spam).
        \item \textbf{Phân loại Đa lớp (Multiclass Classification):} Phân loại đầu vào thành nhiều hơn 2 nhóm (Ví dụ: Nhận dạng chữ số 0-9).
        \item \textbf{Hồi quy Vô hướng (Scalar Regression):} Dự đoán một giá trị liên tục. (Ví dụ: Dự đoán giá nhà, nhiệt độ).
    \end{enumerate}
\end{frame}

\section{2. Phân loại Nhị phân (IMDb Dataset)}

\begin{frame}{Bài toán Đánh giá Cảm xúc (IMDb Dataset)}
    \begin{itemize}
        \item \textbf{Mục tiêu:} Đánh giá xem một bài bình luận phim (movie review) là Tích cực (Positive) hay Tiêu cực (Negative).
        \item \textbf{Dữ liệu:} 50.000 bài đánh giá từ Internet Movie Database, chia đều 25.000 train và 25.000 test. Cân bằng 50\% Positive - 50\% Negative.
        \item \textbf{Tiền xử lý dữ liệu:} Từ ngữ được mã hóa thành các số nguyên (integer), sau đó dùng Multi-hot Encoding để biến đổi thành các vector nhị phân (chỉ chứa 0 và 1) có chiều dài 10.000.
    \end{itemize}
\end{frame}

\begin{frame}{Kiến trúc Mạng Nơ-ron 3 Lớp}
    \begin{columns}
        \column{0.5\textwidth}
        \begin{itemize}
            \item Mạng có 2 lớp ẩn (Hidden Layers), mỗi lớp có 16 nơ-ron (units) và dùng hàm kích hoạt ReLU.
            \item Lớp đầu ra (Output Layer) có 1 nơ-ron, dùng hàm kích hoạt Sigmoid để tính xác suất $[0, 1]$.
            \item Ý nghĩa: Lớp thứ nhất trích xuất đặc trưng cơ bản. Lớp thứ hai phân tích quan hệ phức tạp.
        \end{itemize}
        
        \column{0.5\textwidth}
        \begin{figure}
            \centering
            \includegraphics[width=\textwidth,height=0.7\textheight,keepaspectratio]{../../images/ch04/3_layer_network.cf1b1cd7.png}
            \caption{Cấu trúc mạng 3 lớp cho phân loại IMDb}
        \end{figure}
    \end{columns}
\end{frame}

\begin{frame}{Hàm kích hoạt (Activation): Tại sao lại là ReLU?}
    \begin{columns}
        \column{0.5\textwidth}
        \begin{equation}
            f(x) = \max(0, x)
        \end{equation}
        \begin{itemize}
            \item Không có hàm kích hoạt phi tuyến, mạng nơ-ron nhiều lớp cũng chỉ tương đương 1 lớp tuyến tính (không thể giải bài toán phức tạp).
            \item \textbf{ReLU (Rectified Linear Unit):} Đơn giản là loại bỏ giá trị âm (trả về 0), giữ nguyên giá trị dương. Giúp mạng học rất nhanh và tránh triệt tiêu đạo hàm.
        \end{itemize}
        
        \column{0.5\textwidth}
        \begin{figure}
            \centering
            \includegraphics[width=\textwidth,height=0.7\textheight,keepaspectratio]{../../images/ch04/The-rectified-linear-unit-function.351095bf.png}
            \caption{Đồ thị của hàm ReLU}
        \end{figure}
    \end{columns}
\end{frame}

\begin{frame}{Lớp Output: Hàm Sigmoid}
    \begin{columns}
        \column{0.5\textwidth}
        \begin{equation}
            f(x) = \frac{1}{1 + e^{-x}}
        \end{equation}
        \begin{itemize}
            \item Đầu ra của mạng có thể là số bất kỳ $(-\infty, +\infty)$. Nhưng ta cần xác suất phần trăm (từ 0\% đến 100\%).
            \item \textbf{Sigmoid:} Ép bất kỳ giá trị nào về miền $[0, 1]$.
            \item Kết quả $0.9 \rightarrow 90\%$ khả năng là bình luận Tích cực.
        \end{itemize}
        
        \column{0.5\textwidth}
        \begin{figure}
            \centering
            \includegraphics[width=\textwidth,height=0.7\textheight,keepaspectratio]{../../images/ch04/The-sigmoid-function.eac1368d.png}
            \caption{Đồ thị của hàm Sigmoid}
        \end{figure}
    \end{columns}
\end{frame}

\begin{frame}{Đánh giá Huấn luyện: Biểu đồ Loss}
    \begin{columns}
        \column{0.5\textwidth}
        \begin{itemize}
            \item Dùng \textbf{Binary Crossentropy} làm hàm mất mát.
            \item Đường chấm bi (Training Loss) liên tục giảm $\rightarrow$ Mô hình đang học tốt trên tập huấn luyện.
            \item Đường liền (Validation Loss) giảm rồi đột ngột tăng lại $\rightarrow$ Hiện tượng \textbf{Overfitting (Quá khớp)}. Mô hình "học vẹt" thay vì khái quát hóa.
        \end{itemize}
        
        \column{0.5\textwidth}
        \begin{figure}
            \centering
            \includegraphics[width=\textwidth,height=0.7\textheight,keepaspectratio]{../../images/ch04/imdb_loss_plot.801b28d0.png}
            \caption{Sự thay đổi Loss qua các Epochs}
        \end{figure}
    \end{columns}
\end{frame}

\begin{frame}{Đánh giá Huấn luyện: Biểu đồ Accuracy}
    \begin{columns}
        \column{0.5\textwidth}
        \begin{itemize}
            \item Training Accuracy (chấm bi) tiến sát 100\% (1.0).
            \item Validation Accuracy (đường liền) đạt đỉnh khoảng 88\% tại epoch thứ 4, sau đó giảm dần.
            \item \textbf{Giải pháp:} Nên dừng huấn luyện sớm (Early Stopping) tại epoch thứ 4 để lấy mô hình tốt nhất, tránh Overfitting.
        \end{itemize}
        
        \column{0.5\textwidth}
        \begin{figure}
            \centering
            \includegraphics[width=\textwidth,height=0.7\textheight,keepaspectratio]{../../images/ch04/imdb_accuracy_plot.bf0cb7ef.png}
            \caption{Sự biến thiên của Accuracy}
        \end{figure}
    \end{columns}
\end{frame}

\section{3. Phân loại Đa lớp (Reuters Dataset)}

\begin{frame}{Bài toán Reuters (Multiclass Classification)}
    \begin{itemize}
        \item \textbf{Mục tiêu:} Phân loại các tin tức của Reuters vào đúng 46 chủ đề khác nhau (Topics).
        \item Đây là bài toán \textbf{Single-label, Multiclass} (Mỗi bài báo chỉ thuộc đúng 1 chủ đề, nhưng có tới 46 chủ đề để chọn).
        \item \textbf{Kiến trúc:} Do số lượng phân lớp lớn (46), không thể dùng lớp ẩn có 16 nơ-ron (bị "thắt cổ chai" thông tin). Phải dùng 64 nơ-ron cho mỗi lớp.
        \item Lớp cuối cùng: 46 nơ-ron kết hợp hàm kích hoạt \textbf{Softmax} (tạo ra phân phối xác suất có tổng bằng 1).
    \end{itemize}
\end{frame}

\begin{frame}{Hàm mất mát \& Biểu đồ Loss (Reuters)}
    \begin{columns}
        \column{0.5\textwidth}
        \begin{itemize}
            \item Dùng \textbf{Categorical Crossentropy}: Chuyên dụng đo lường khoảng cách giữa hai phân phối xác suất đa chiều.
            \item Hiện tượng Overfitting cũng xảy ra nhanh chóng sau epoch thứ 9. Loss trên tập validation bắt đầu tăng mạnh.
        \end{itemize}
        
        \column{0.5\textwidth}
        \begin{figure}
            \centering
            \includegraphics[width=\textwidth,height=0.7\textheight,keepaspectratio]{../../images/ch04/reuters_loss_plot.6e487e1a.png}
            \caption{Training và Validation Loss (Reuters)}
        \end{figure}
    \end{columns}
\end{frame}

\begin{frame}{Đánh giá Accuracy}
    \begin{columns}
        \column{0.5\textwidth}
        \begin{itemize}
            \item Accuracy trên tập Validation đạt cao nhất khoảng ~80\% trước khi quá khớp.
            \item So với xác suất chọn ngẫu nhiên (chỉ khoảng 19\% nếu đoán bừa theo tần suất), mức 80\% là một kết quả tuyệt vời.
        \end{itemize}
        
        \column{0.5\textwidth}
        \begin{figure}
            \centering
            \includegraphics[width=\textwidth,height=0.7\textheight,keepaspectratio]{../../images/ch04/reuters_accuracy_plot.b74dee12.png}
            \caption{Training và Validation Accuracy}
        \end{figure}
    \end{columns}
\end{frame}

\begin{frame}{Nhưng 80\% đã là con số thực? (Top-3 Accuracy)}
    \begin{columns}
        \column{0.5\textwidth}
        \begin{itemize}
            \item Đôi khi một bài báo có nội dung đan xen nhiều mảng. Việc mô hình dự đoán nhãn đúng nằm trong Top-3 lựa chọn cao nhất cũng là rất giá trị.
            \item \textbf{Top-3 Accuracy:} Đo lường tỷ lệ nhãn thật (Ground-truth) nằm trong 3 nhãn có xác suất cao nhất của mô hình dự đoán.
            \item Chỉ số này có thể lên tới >90\%.
        \end{itemize}
        
        \column{0.5\textwidth}
        \begin{figure}
            \centering
            \includegraphics[width=\textwidth,height=0.7\textheight,keepaspectratio]{../../images/ch04/reuters_top_3_accuracy_plot.a9e13ec0.png}
            \caption{Chỉ số Top-3 Accuracy}
        \end{figure}
    \end{columns}
\end{frame}

\section{4. Hồi quy Vô hướng (Boston Housing)}

\begin{frame}{Bài toán Dự đoán Giá nhà Boston (Scalar Regression)}
    \begin{itemize}
        \item \textbf{Mục tiêu:} Dự đoán giá nhà trung bình tại vùng ngoại ô Boston (thập niên 1970) dựa trên 13 đặc trưng (tỷ lệ tội phạm, thuế, số phòng...).
        \item \textbf{Khác biệt lớn:} Đầu ra không phải là "xác suất" hay "lớp", mà là \textbf{một con số liên tục} (ví dụ: \$25.000). 
        \item Lớp Output: Chỉ 1 nơ-ron và \textbf{không dùng hàm kích hoạt} (Linear Layer) để không giới hạn miền giá trị dự đoán.
        \item \textbf{Hàm mất mát:} Dùng \textbf{MSE} (Mean Squared Error).
        \item \textbf{Độ đo (Metric):} Dùng \textbf{MAE} (Mean Absolute Error - Sai số tuyệt đối trung bình). Giúp hiểu dễ hơn: MAE=0.5 tức là lệch \$500.
    \end{itemize}
\end{frame}

\begin{frame}{Đặc thù Dữ liệu Nhỏ \& K-Fold Cross-Validation}
    \begin{columns}
        \column{0.5\textwidth}
        \begin{itemize}
            \item Tập dữ liệu chỉ có 506 mẫu $\rightarrow$ Nếu lấy 100 mẫu làm tập Validation, sai số đánh giá sẽ dao động rất lớn, phụ thuộc vào việc "bốc trúng" dữ liệu nào.
            \item \textbf{Giải pháp: K-Fold Cross Validation} (Thẩm định chéo K-lần).
            \item Chia dữ liệu thành K phần (ví dụ $K=4$). Lặp lại 4 lần: 1 phần làm Validation, 3 phần làm Training.
            \item Điểm số cuối cùng = Trung bình của 4 lần huấn luyện.
        \end{itemize}
        
        \column{0.5\textwidth}
        \begin{figure}
            \centering
            \includegraphics[width=\textwidth,height=0.7\textheight,keepaspectratio]{../../images/ch04/3-fold-cross-validation.40bb5356.png}
            \caption{Mô phỏng 3-Fold Cross-Validation}
        \end{figure}
    \end{columns}
\end{frame}

\begin{frame}{Đánh giá Hiệu suất Hồi quy (MAE Plot)}
    \begin{columns}
        \column{0.5\textwidth}
        \begin{itemize}
            \item Khác với Accuracy (càng cao càng tốt), MAE (sai số) càng thấp càng tốt.
            \item Trực quan hóa giá trị trung bình MAE qua K-folds trên 500 epochs.
            \item Đường đồ thị có độ rung (nhiễu) cao ở những epochs đầu.
        \end{itemize}
        
        \column{0.5\textwidth}
        \begin{figure}
            \centering
            \includegraphics[width=\textwidth,height=0.7\textheight,keepaspectratio]{../../images/ch04/california_housing_validation_mae_plot.af306c57.png}
            \caption{Validation MAE qua 500 epochs}
        \end{figure}
    \end{columns}
\end{frame}

\begin{frame}{Phân tích Chuyên sâu (MAE Zoomed Plot)}
    \begin{columns}
        \column{0.5\textwidth}
        \begin{itemize}
            \item Bỏ qua 10 epochs đầu tiên (do sai số quá lớn làm lệch thang biểu đồ).
            \item Vẽ đường làm trơn bằng trung bình động hàm mũ (Exponential Moving Average) để quan sát rõ hơn xu hướng.
            \item Sai số chạm đáy ở khoảng 120-140 epochs, sau đó bắt đầu tăng vọt $\rightarrow$ Overfitting trên tập dữ liệu hồi quy.
        \end{itemize}
        
        \column{0.5\textwidth}
        \begin{figure}
            \centering
            \includegraphics[width=\textwidth,height=0.7\textheight,keepaspectratio]{../../images/ch04/california_housing_validation_mae_plot_zoomed.928f390d.png}
            \caption{Đồ thị làm trơn (Zoomed)}
        \end{figure}
    \end{columns}
\end{frame}

\section{Tổng kết}

\begin{frame}{Tổng kết Chương 4}
    \begin{itemize}
        \item \textbf{Phân loại Nhị phân:} Dùng 1 nơ-ron output với hàm \textbf{Sigmoid} và Loss \textbf{Binary Crossentropy}.
        \item \textbf{Phân loại Đa lớp (Single-label):} Dùng N nơ-ron với hàm \textbf{Softmax} và Loss \textbf{Categorical Crossentropy}. Cần chú ý nút thắt cổ chai kiến trúc.
        \item \textbf{Hồi quy (Regression):} Không dùng hàm kích hoạt ở output layer. Loss phổ biến là \textbf{MSE}, và metric theo dõi là \textbf{MAE}.
        \item \textbf{Dữ liệu nhỏ:} Bắt buộc sử dụng K-Fold Cross Validation.
        \item \textbf{Kẻ thù chung:} Hiện tượng quá khớp (Overfitting) xảy ra ở tất cả các bài toán. Việc theo dõi biểu đồ Loss/Metric là bắt buộc.
    \end{itemize}
\end{frame}

\end{document}
